\documentclass[10pt,a4paper]{amsart}
\usepackage[margin=19mm]{geometry}
\usepackage{amsmath,amssymb,mathtools}
\usepackage[scr=rsfs,cal=euler]{mathalfa}
\usepackage{xfrac}
\usepackage[unicode,hidelinks,bookmarks=false]{hyperref}
\newcommand{\ie}{i.e.\ }
\newcommand{\cf}{cf.\ }
\newcommand{\resp}{respectively\ }
\newcommand{\Subsection}{\subsection}
\providecommand{\nobreakdash}{\nobreak\hbox{-}}
\setlength{\emergencystretch}{2em}
\multlinegap0pt
\usepackage{mathtools}
\usepackage{amssymb}
\usepackage{xcolor}
\usepackage{tikz}
\newtheorem{theorem}{Theorem}
\newtheorem{lemma}{Lemma}
\title{Caged subsequences in permutations}
\author{Niranjan Balachandran, Omkar Ramdas, Umesh Shankar}
\date{Source: arXiv:2608.09770v2 (CC BY 4.0)}
\begin{document}
\maketitle

\noindent\emph{Source and licence.} Caged subsequences in permutations. Version arXiv:2608.09770v2 (CC BY 4.0), distributed by arXiv under the Creative Commons Attribution 4.0 International licence, http://creativecommons.org/licenses/by/4.0/, as recorded in the arXiv licence metadata for this exact version and shown on the abstract page https://arxiv.org/abs/2608.09770v2. Full licence notice: \texttt{licenses/arxiv-2608.09770-CC-BY-4.0.txt}.\par\smallskip

\noindent\textit{Editorial scope.} Counted unit: the article's theorem expectation (source0.tex lines 207--220). The article's complete original proof of it is delivered with it (source0.tex lines 656--729, one \texttt{proof} environment of 74 lines, which the article places away from the statement). No mathematics is written, repaired or re-derived: the statement and the proof are the article's own lines, in the article's own notation, and no retained line is substituted or re-flowed.  No further source-local context is needed: every cross-reference of every retained line resolves inside the two delivered spans.  Nothing was omitted from any retained span, so the package carries no omission record.  No retained line contains a citation, so the wrapper carries no bibliography and no bibliography entry.  No retained line contains a figure environment, an image-inclusion command or drawing code, so no image is delivered and the package declares no asset dependency.  Editorial formatting only, in addition: an AMS \texttt{amsart} preamble that restates the article's own declarations and macros used by the retained lines, namely the environments \texttt{theorem}, \texttt{lemma} on separate counters, together with the standard packages those lines need.  The retained environments print in this document as Theorem 1, Lemma 1 in order of appearance, whereas the article prints the same items with the numbering of its own full document.  The complete native source of the article as distributed in the arXiv e-print for this exact version is bundled unchanged as \texttt{sources/207271/source0.tex}.
\begin{theorem}\label{expectation}
For any $\omega(n) \rightarrow \infty$ as $n \rightarrow \infty$,
\[
  \lim_{n \rightarrow \infty} \mathbb{P}\left( c(\pi) \geq \left(1-\frac{\omega(n)}{\sqrt{n}}\right)n \right) = 1.
\]
Also, 
\[
  \lim_{n \rightarrow \infty} \mathbb{P}\left( c(\pi) \leq \left(1-\frac{1}{\omega(n)\sqrt{n}}\right)n \right) = 1.
\]
Consequently, 
\[
\mathbb{E}[c(\pi)] = (1 - o(1))n .
\]
\end{theorem}

\begin{proof}[Proof of theorem\ref{expectation}]
 Note that a uniformly random permutation can be generated by choosing the $X_i \sim U[0,1]$ i.i.d. for $1 \leq i \leq n$. So, we can work on these $X_i$'s instead. A long caged sequence will appear if two extreme $X_i$'s(say $X_{i_1}$ and $X_{i_2} $) assume the extreme values in the interval $[0, 1]$ and many other $X_j$'s whose indices fall between $i_1$ and $i_2$ take the values in between. We will now implement this idea mathematically. \\
 Without loss of generality, let $\omega(n)$ be any small function $\left(o(\sqrt{n})\right)$ such that $\omega(n) \xrightarrow{n \rightarrow \infty } \infty$. Take $\displaystyle\epsilon = \frac{\omega(n)}{\sqrt{n}}$. Define $\delta := \displaystyle \frac{1 - (1 - \epsilon)^{\frac{1}{3}}}{2} $. Observe that $0<\delta < < 1$ for $n$ is large. Let $\mathcal{E}_1 $ be the event that for some $i \in [1, \delta n]$, $X_i \in [0, \delta ]$. Similarly, let $\mathcal{E}_2 $ be the event that for some $j \in [(1 -\delta) n, n]$, $X_j \in [(1 - \delta), 1]$. Then 
 \[
 \mathbb{P}\left( \mathcal{E}_1 \wedge \mathcal{E}_2 \right) = 1 - \mathbb{P}\left( \bar{\mathcal{E}}_1 \vee \bar{\mathcal{E}}_2 \right)
 \] 
 Since the random variables $X_i$ are independent, $ \mathbb{P}\left( \bar{\mathcal{E}}_1 \right) = \mathbb{P} \left( \bar{\mathcal{E}}_2 \right) = (1 - \delta )^{\delta n} $. By union bound,
 \[
 \mathbb{P}\left( \bar{\mathcal{E}}_1 \vee \bar{\mathcal{E}}_2 \right) \leq 2(1 - \delta )^{\delta n} \leq 2e^{-\delta^2 n}.
 \] 
 So, 
 \[
 \mathbb{P}\left( \mathcal{E}_1 \wedge \mathcal{E}_2 \right) \geq 1 -  2e^{-\delta^2 n}.
 \]
 Define $Y : = \displaystyle \sum_{k \in (\delta n, (1 - \delta)n)} \displaystyle  \mathbb{1}_{X_k \in (\delta , (1- \delta ))} $. Observe that $Y$ corresponds to the length of the caged sequence. Using linearity of expectation, $\mu := \mathbb{E}( Y) = (1 - 2\delta)^2 n$. 
 Let $\mathcal{E}_3$ be the event that $Y$ is at least $(1 - \gamma) \mu$ for some small $\gamma >0$. Using the Chernoff bound for the sum of Bernoulli i.i.ds,
 \[
 \mathbb{P}\left( \bar{\mathcal{E}}_3 \right) = \mathbb{P}\left( Y \leq (1 - \gamma)\mu  \right) \leq \displaystyle e^{-\gamma^2 \mu / 2}.
 \]
 Finally, we use independence to get  
 \[
 \mathbb{P}\left( \mathcal{E}_1 \wedge \mathcal{E}_2 \wedge \mathcal{E}_3  \right) \geq \left( 1 -  2e^{-\delta^2 n} \right) \left( 1 - e^{-\gamma^2 \mu / 2} \right).
 \]
 Put $\gamma = 2 \delta$. Then the above statement implies $ \mathbb{P}\left(c(\pi) \geq (1 - 2 \delta)^3 n \right) \longrightarrow 1$ as $ n \rightarrow \infty$. \\
 By plugging the values of $\delta$ and $ \epsilon $, $\mathbb{P}\left( c(\pi) \geq \left( 1- \frac{\omega(n)}{\sqrt{n}} \right)n \right) \longrightarrow 1 $ as $n$ tends to infinity. \\
 To prove the second statement of the theorem, we need the following lemma:
\begin{lemma}\label{lemma 4.2}
   Let \(X,Y\sim U[0,1]\) be i.i.d. Let \(D = |X-Y|\) and \(N\mid D \sim \text{Bin}(n,D)\). Then for any $d = d(n)$ where $d=o(n)$, we have
   $$ \mathbb{P}(N>n-d(n)) \leq \exp{\left(-\frac{9}{8}\cdot d(n)\right)} + 16\left(\frac{d(n)}{n}\right)^2.$$
\end{lemma}
\begin{proof}
Set $t = n - d(n)$ with $d(n) = o(n)$. Write $M = n-N, Q= 1-D$, and define $M$ conditioned on $D$ to be  $\text{Bin}\big(n,Q\big)$. It is straightforward to check that $Q$ has  density $f_Q(q) = 2q$ for $q \in [0, 1]$.\\
Now,
\begin{align*}
    \mathbb{P}(N>n-d(n)) & = \mathbb{P}\left( M < d(n) \right) \\
   &= \mathbb{E}_Q\left[ \mathbb{P}(M<d(n) ) \mid Q\right] \\
   & = \int_0^1 \mathbb{P}\big(M<d(n)\mid Q=q\big)\ 2q\ dq.
\end{align*}
 For $0\le q\le q_0:=4d(n)/n$, we have $\mathbb{P}(M<d(n)\mid q)\le 1$, so
    \begin{align}\label{eq:9}
        \int_0^{q_0} \mathbb{P}(M<d(n)\mid q)\ 2q\  dq
    \le \int_0^{q_0} 2q\ dq
    = 16\left(\frac{d(n)}{n}\right)^2.
    \end{align}
    For $q_0 < q \leq 1$, the mean of $M$ is $nq > 4d(n)$, and as $M<d(n)$ is a lower tail event, a Chernoff-type bound gives 
    \begin{align*}
         \mathbb{P}(M<d(n)\mid q) & = \mathbb{P}\left(  M < \mathbb{E}(M)  -(nq - d(n))\right) \\
        & \leq \exp{\left( -\frac{(nq - d(n))^2}{2nq} \right)} \\
       & \leq \exp{\left( -\frac{9}{8}\cdot d(n) \right)},
    \end{align*}
    for all $q_0 < q \leq 1$. Thus,
    \begin{align}
        \int_{q_0}^1 \mathbb{P}(M<d(n)\mid q)\ 2q\  dq \le \exp\left(-\frac{9}{8}\cdot d(n)\right).
    \end{align}
Thus, we have 
$$
\mathbb{P}(N>n-d(n)) \leq 16\left(\frac{d(n)}{n}\right)^2 + \exp\left(-\frac{9}{8}\cdot d(n)\right).
$$
\end{proof}
We are now in a position to prove the second statement of theorem \ref{expectation}. \\
Set $d(n)=o(n)$. We want to estimate the probability that there exists any caged sequence of length $n - d(n)$. If this holds, there must exist a pair $i, j $ with $j - i  \geq n - d(n) - 1$ and $c[i, j] \geq n - d(n)$. Note that there are at most $(d(n))^2$ pairs with $j - i \geq n - d(n) - 1$. Pick any such pair $(i, j)$. In terms of earlier notation, we have the following.
\[
   c[i, j] = 2 \, +  \displaystyle \sum_{i < k < j} \displaystyle  \mathbb{1}_{X_i < X_k <X_j \text{ or } X_i > X_k > X_j}. 
\]
Write $N = c[i, j] - 2$. Clearly, $N \mid D \sim \text{Bin}(j - i - 1, D)$ where $D = |X_i - X_j|$. Now, $j - i - 1 \geq n - d(n) - 2 = (1-o(1))n$. Applying the lemma \ref{lemma 4.2}, 
\[
  \mathbb{P}\left( c[i, j] \geq n - d(n) \right) \leq  \exp{\left(-\frac{9}{8}\cdot (d(n) + 2) \right)} + 16\left(\frac{d(n) + 2}{(1+o(1))n}\right)^2.
\]
Using the union bound, 
\[
\mathbb{P}\left( \text{ There exists such a pair $i, j$ with } c[i, j] \geq n - d(n) \right) \leq \frac{d(n)^2}{\exp{\left(\frac{9}{8}\cdot (d(n) + 2) \right)}} + 16\frac{(d(n) + 2)^4}{(1+o(1))n^2}.
\]
For $d(n) = \frac{\sqrt{n}}{\omega(n)}$ with $\omega(n) = o(\sqrt{n})$ and $\omega(n) \rightarrow \infty $, the RHS in the above expression tends to $0$; this completes the proof.
\end{proof}

\end{document}
